\begin{thebibliography}{33}
\bibitem{r1} M. Ajtai, J. Komlos, and E. Szemeredi. Deterministic simulation in LOGSPACE. In Proceedings of the Nineteenth Annual ACM Symposium on Theory of Computing (STOC), pages 132–140, New York, NY, USA, 1987. ACM. doi:10.1145/28395.28410.
\bibitem{r2} Miklos Ajtai and Avi Wigderson. Deterministic simulation of probabilistic constant depth circuits. Advances in Computing Research, 5(199-222):1, 1989.
\bibitem{r3} Noga Alon, Oded Goldreich, Johan Håstad, and René Peralta. Simple constructions of almost k-wise independent random variables. Random Structures \& Algorithms, 3(3):289–304, 1992.
\bibitem{r4} Louay Bazzi and Nagi Nahas. Small-bias is not enough to hit read-once CNF. Theory of Computing Systems, 60(2):324–345, February 2017. doi:10.1007/s00224-016-9680-6.
\bibitem{r5} Mark Braverman, Gil Cohen, and Sumegha Garg. Pseudorandom pseudo-distributions with near-optimal error for read-once branching programs. SIAM Journal on Computing, 0(0):STOC18–242–STOC18–299, 2020. doi:10.1137/18M1197734.
\bibitem{r6} Mark Braverman, Anup Rao, Ran Raz, and Amir Yehudayoff. Pseudorandom generators for regular branching programs. SIAM Journal on Computing, 43(3):973–986, 2014.
\bibitem{r7} J. Brody and E. Verbin. The coin problem and pseudorandomness for branching programs. In 2010 IEEE 51st Annual Symposium on Foundations of Computer Science (FOCS), pages 30–39, October 2010. doi:10.1109/FOCS.2010.10.
\bibitem{r8} Suresh Chari, Pankaj Rohatgi, and Aravind Srinivasan. Improved algorithms via approximations of probability distributions. Journal of Computer and System Sciences, 61(1):81–107, 2000. doi:10.1006/jcss.1999.1695.
\bibitem{r9} Eshan Chattopadhyay, Pooya Hatami, Omer Reingold, and Avishay Tal. Improved pseudorandomness for unordered branching programs through local monotonicity. In Proceedings of the 50th Annual ACM Symposium on Theory of Computing (STOC), pages 363–375, New York, NY, USA, 2018. ACM. doi:10.1145/3188745.3188800.
\bibitem{r10} Sitan Chen, Thomas Steinke, and Salil Vadhan. Pseudorandomness for read-once, constantdepth circuits. arXiv preprint, 2015. arXiv:1504.04675.
\bibitem{r11} Anindya De. Pseudorandomness for permutation and regular branching programs. In Proceedings of the 26th Annual IEEE 26th Annual Conference on Computational Complexity (CCC), pages 221–231. IEEE, 2011.
\bibitem{r12} Anindya De, Omid Etesami, Luca Trevisan, and Madhur Tulsiani. Improved pseudorandom generators for depth 2 circuits. In Approximation, randomization, and combinatorial optimization, volume 6302 of Lecture Notes in Computer Science, pages 504–517. Springer, Berlin, 2010. doi:10.1007/978-3-642-15369-3\_38.
\bibitem{r13} Dean Doron, Pooya Hatami, and William M Hoza. Near-optimal pseudorandom generators for constant-depth read-once formulas. In 34th Computational Complexity Conference (CCC), 2019.
\bibitem{r14} Michael A. Forbes and Zander Kelley. Pseudorandom generators for read-once branching programs, in any order. In Proceedings of the 59th Annual IEEE Symposium on Foundations of Computer Science (FOCS). IEEE, 2018.
\bibitem{r15} Anat Ganor and Ran Raz. Space pseudorandom generators by communication complexity lower bounds. In LIPIcs-Leibniz International Proceedings in Informatics, volume 28. Schloss Dagstuhl-Leibniz-Zentrum fuer Informatik, 2014.
\bibitem{r16} Parikshit Gopalan, Raghu Meka, Omer Reingold, Luca Trevisan, and Salil Vadhan. Better pseudorandom generators from milder pseudorandom restrictions. In Proceedings of the 53rd Annual IEEE Symposium on Foundations of Computer Science (FOCS), pages 120–129. IEEE, 2012.
\bibitem{r17} Parikshit Gopalan and Amir Yehudayoff. Inequalities and tail bounds for elementary symmetric polynomial with applications. arXiv preprint, 2014. arXiv:1402.3543.
\bibitem{r18} Elad Haramaty, Chin Ho Lee, and Emanuele Viola. Bounded independence plus noise fools products. SIAM Journal on Computing, 47(2):493–523, 2018. doi:10.1137/17M1129088.
\bibitem{r19} William M. Hoza and David Zuckerman. Simple optimal hitting sets for small-success RL. In Proceedings of the 59th Annual IEEE Symposium on Foundations of Computer Science (FOCS). IEEE, 2018.
\bibitem{r20} Russell Impagliazzo, Noam Nisan, and Avi Wigderson. Pseudorandomness for network algorithms. In Proceedings of the Twenty-Sixth Annual ACM Symposium on Theory of Computing (STOC), pages 356–364. ACM, 1994.
\bibitem{r21} Michal Koucký, Prajakta Nimbhorkar, and Pavel Pudlák. Pseudorandom generators for group products. In Proceedings of the 43rd Annual ACM Symposium on Theory of Computing (STOC), pages 263–272. ACM, New York, 2011. doi:10.1145/1993636.1993672.
\bibitem{r22} Chin Ho Lee. Fourier bounds and pseudorandom generators for product tests. In Amir Shpilka, editor, 34th Computational Complexity Conference (CCC), volume 137 of Leibniz International Proceedings in Informatics (LIPIcs), pages 7:1–7:25, Dagstuhl, Germany, 2019. Schloss Dagstuhl–Leibniz-Zentrum fuer Informatik. doi:10.4230/LIPIcs.CCC.2019.7.
\bibitem{r23} Chin Ho Lee and Emanuele Viola. More on bounded independence plus noise: Pseudorandom generators for read-once polynomials. In Electronic Colloquium on Computational Complexity (ECCC), volume 24, page 167, 2017.
\bibitem{r24} Raghu Meka, Omer Reingold, and Avishay Tal. Pseudorandom generators for width-3 branching programs. In Proceedings of the 51st Annual ACM Symposium on Theory of Computing (STOC), pages 626–637. ACM, New York, 2019.
\bibitem{r25} Joseph Naor and Moni Naor. Small-bias probability spaces: Efficient constructions and applications. SIAM Journal on Computing, 22(4):838–856, 1993.
\bibitem{r26} Noam Nisan. Pseudorandom generators for space-bounded computation. Combinatorica, 12(4):449–461, 1992.
\bibitem{r27} Noam Nisan. $\mathrm{RL}\subseteq\mathrm{SC}$. computational complexity, 4(1):1–11, March 1994. doi:10.1007/ BF01205052.
\bibitem{r28} Noam Nisan and David Zuckerman. Randomness is linear in space. Journal of Computer and System Sciences, 52(1):43–52, 1996. doi:10.1006/jcss.1996.0004.
\bibitem{r29} Omer Reingold. Undirected connectivity in log-space. Journal of the ACM, 55(4):Art. 17, 24, 2008. doi:10.1145/1391289.1391291.
\bibitem{r30} Omer Reingold, Thomas Steinke, and Salil Vadhan. Pseudorandomness for regular branching programs via Fourier analysis. In Approximation, Randomization, and Combinatorial Optimization. Algorithms and Techniques, pages 655–670. Springer, 2013.
\bibitem{r31} Michael E. Saks and Shiyu Zhou. $\mathrm{BPHSPACE}(S)\subseteq\mathrm{DSPACE}(S^{3/2})$. Journal of Computer and System Sciences, 58(2):376–403, 1999.
\bibitem{r32} Thomas Steinke. Pseudorandomness for permutation branching programs without the group theory. In Electronic Colloquium on Computational Complexity (ECCC), 2012.
\bibitem{r33} Thomas Steinke, Salil Vadhan, and Andrew Wan. Pseudorandomness and Fourier-growth bounds for width-3 branching programs. Theory of Computing, 13(12):1–50, 2017. doi: 10.4086/toc.2017.v013a012.
\end{thebibliography}
